% Claim proof file for row claim_3_19 -- "MarginalizingOutThe".
%
% Proof lifted verbatim from lecture-notes/lecture_notes/graphs.tex lines
% 1177--1210 (the LN's `\Claude{...}` wrapper stripped). The LN proof
% already follows the three-section Node-sets / Directed-edges /
% Bidirected-edges skeleton called for in the worker prompt; no
% adaptation was needed.
%
% Compiles standalone via the `subfiles` package.

\documentclass[main]{subfiles}

\begin{document}
% ref: claim_3_19 (proof, with statement restated for readability)
% title: MarginalizingOutThe
\def\rowref{claim\_3\_19}
\phantomsection\label{claim_3_19}
%
% Cross-references: see claim_<ref>_statement_<title>.tex in this folder
% for the corresponding statement.

% --- statement (restated) ---------------------------------------------------
\begin{Lem}[Marginalizing out the output part of splitted nodes equals hard intervention]
     Let $G=(J,V,E,L)$ be a CDMG and $W \ins V$ be a subset of output nodes from $G$.
     Then the CDMG that arises by first splitting the nodes on $W$
      and then marginalizing out the nodes from $W^o$
     can be identified with the CDMG that arises by hard intervention on $W$:
     \[ G_{\doit(W)} \cong \lp G_{\swig(W)}\rp^{\sm W^o}, \qquad w \mapsto w^i. \]
\end{Lem}

% --- proof ------------------------------------------------------------------
\begin{proof}
We establish the isomorphism $G_{\doit(W)} \cong \lp G_{\swig(W)}\rp^{\sm W^o}$ via the map $\phi$ that is the identity on $J \cup (V \sm W)$ and sends $w^i \mapsto w$ for each $w \in W$.

Recall:
\begin{itemize}
    \item $G_{\doit(W)}$ has input nodes $J \cup W$, output nodes $V \sm W$, directed edges $E \sm \{v \tuh w \mid w \in W\}$, and bidirected edges $L \sm \{v \huh w \mid w \in W\}$.
    \item $G_{\swig(W)}$ has input nodes $J \dcup W^i$, output nodes $(V \sm W) \dcup W^o$, directed edges $\{v_1^i \tuh v_2^o \mid v_1 \tuh v_2 \in E\}$, and bidirected edges $\{v_1^o \huh v_2^o \mid v_1 \huh v_2 \in L\}$ (where $v^i = v^o = v$ for $v \notin W$).
\end{itemize}

\medskip\noindent\emph{Node sets.}
After marginalizing out $W^o$:
input nodes $J \dcup W^i$ (unchanged) and output nodes $(V \sm W) \dcup W^o \sm W^o = V \sm W$.
Under $\phi$ ($w^i \mapsto w$) this matches $G_{\doit(W)}$: input nodes $J \cup W$, output nodes $V \sm W$.

\medskip\noindent\emph{Directed edges.}
In $G_{\swig(W)}$, nodes $w^o \in W^o$ have no outgoing directed edges (all outgoing edges of $w$ go to $w^i$, not $w^o$).
Therefore no directed walk can pass through a node in $W^o$ as an intermediate node.
Hence the marginalization step does not create any new directed edges: $\ul{v} \tuh \ol{v} \in (E_{\swig(W)})^{\sm W^o}$ iff $\ul{v} \tuh \ol{v} \in E_{\swig(W)}$ with $\ul{v},\ol{v} \notin W^o$.

Such an edge corresponds to some $v_1 \tuh v_2 \in E$ with $\ul{v} = v_1^i$ and $\ol{v} = v_2^o$.
Since $\ol{v} \notin W^o$, we have $v_2 \notin W$, so $v_2^o = v_2$.
Under $\phi$: $\phi(\ul{v}) \tuh \phi(\ol{v}) = v_1 \tuh v_2 \in E$ with $v_2 \notin W$,
i.e.\ $v_1 \tuh v_2 \in E_{\doit(W)}$.
Conversely, every edge in $E_{\doit(W)}$ lifts to an edge in $E_{\swig(W)}$ with both endpoints outside $W^o$.

\medskip\noindent\emph{Bidirected edges.}
Let $\ul{v}, \ol{v} \in V \sm W$, $\ul{v} \neq \ol{v}$.
Since nodes in $W^o$ have no outgoing directed edges, the only bifurcations in $G_{\swig(W)}$ with intermediate nodes in $W^o$ reduce to a single bidirected edge (the directed arms of the bifurcation cannot extend through $W^o$).
So $\ul{v} \huh \ol{v} \in (L_{\swig(W)})^{\sm W^o}$ iff $\ul{v} \huh \ol{v} \in L_{\swig(W)}$, i.e.\ iff $\ul{v} \huh \ol{v} \in L$ (since $\ul{v}^o = \ul{v}$ and $\ol{v}^o = \ol{v}$ for $\ul{v}, \ol{v} \notin W$).
Under $\phi$ this matches $\ul{v} \huh \ol{v} \in L_{\doit(W)} = L \sm \{v \huh w \mid w \in W\}$, which equals $\ul{v} \huh \ol{v} \in L$ since $\ul{v}, \ol{v} \notin W$.

\medskip
Since $\phi$ provides a bijection on nodes that preserves input/output status and all edges, the two CDMGs are isomorphic.
\end{proof}

\end{document}
